% !TEX root = ../../main.tex
\section{Bố cục của báo cáo đồ án tốt nghiệp}
\label{sec:bo_cuc}
Báo cáo đồ án tốt nghiệp được tổ chức thành 4 chương chính và phần kết luận, tài liệu tham khảo:
\begin{itemize}
    \item \textbf{Chương 1: Tổng quan về đề tài} -- Trình bày bối cảnh ra đời, tính cấp thiết của đề tài, các mục tiêu nghiên cứu tổng quát và cụ thể, nội dung nghiên cứu, phương pháp tiếp cận, phạm vi nghiên cứu và bố cục của báo cáo.
    \item \textbf{Chương 2: Cơ sở lý thuyết và công nghệ} -- Trình bày chi tiết cơ sở lý thuyết về Kiến trúc Lục giác, công nghệ tìm kiếm lai Agentic Hybrid RAG \& RRF, đồ thị tri thức GraphRAG 4 thực thể, tác nhân LangGraph ReAct trong môi trường E2B Sandbox, kỹ thuật CAS \& BYOK bảo mật, đồng thời đưa ra các luận giải khoa học khẳng định lý do lựa chọn các công nghệ này thay vì các giải pháp thay thế.
    \item \textbf{Chương 3: Phân tích và thiết kế hệ thống} -- Khảo sát các luồng nghiệp vụ học thuật, đặc tả các ca sử dụng (Use Case) chuẩn BA, thiết kế biểu đồ tuần tự UML, kiến trúc Monorepo 3 tầng và thiết kế chi tiết 16 bảng cơ sở dữ liệu vật lý kèm từ điển dữ liệu.
    \item \textbf{Chương 4: Cài đặt và kết quả thực nghiệm} -- Trình bày cấu hình môi trường triển khai thực tế, giao diện trực quan của các phân hệ chức năng chính, kết quả kiểm thử tự động 100\% qua 4 Cổng xác minh, cùng các số liệu đo lường định lượng thực tế từ hệ thống và đối chiếu trên nền tảng Langfuse Cloud.
    \item \textbf{Kết luận và hướng phát triển} -- Tổng kết các kết quả đạt được của đề tài, chỉ ra các hạn chế còn tồn tại và đề xuất các hướng phát triển hoàn thiện hệ thống trong tương lai.
\end{itemize}
